\section{The information in the endpoint}\label{sec:fisher}

This note asks a statistical question about Kagey's board. If we only see the
bin where the ball lands, how much do we learn
about the switching probability? The answer is that the bin keeps about a
fraction $1/(2T)$ of the Fisher information of the whole path, so about
$1/(2\log N)$ at Kagey's crossing. We use the notation of
Section~\ref{sec:setting}. The first step is left or right with probability
$\frac12$ each, and each later step repeats the one before it with probability
$p=1-q$, where $0<q<1$, and $T=Nq$.
Put $\rho=1-2q$, $u=q/p$ and $y=Nu=T/p$. Let $J$ be the number of
switches, $K$ the number of right steps, $X=2K-N$ the endpoint and $W=X/N$. A path with $J$ switches has probability
$\frac12p^{N-1-J}q^J$. So $J\sim\operatorname{Bin}(N-1,q)$, and the path
carries the Fisher information $(N-1)/(pq)$ about $q$. The endpoint carries
$I_K(q)=\sum_{k=0}^N(\partial_q\Prob(K=k))^2/\Prob(K=k)$, and we call
$e_N(q)=I_K(q)\,pq/(N-1)$ the \emph{efficiency} of the endpoint. It is the
fraction of the information left when we see only the endpoint.

The endpoint sees the switches only through how far the walk spreads. A
normal $X$ with mean $0$ and variance $Np/q$, close to $\Var X$ for large
$T$, would carry the information $\frac12(\partial_q\log(Np/q))^2=1/(2p^2q^2)$,
which is the fraction $1/(2(N-1)pq)\approx1/(2T)$. Theorem~\ref{thm:fisher}
confirms this.

\begin{proposition}[a correlation ratio]\label{prop:eta}
For $N\ge2$ and $0<q<1$,
\[
 e_N(q)=\frac{\Var(\E[J\mid K])}{\Var J}.
\]
For $0<k<N$ we have $\E[J\mid K=k]=\frac{d}{dt}\log S_{k,N-k}(e^t)$ at
$t=\log(q/p)$, where $S_{a,b}$ is the polynomial of
Section~\ref{sec:setting}. For $k\in\{0,N\}$ we have $J=0$.
\end{proposition}

\begin{proof}
The score of a path is $(J-(N-1)q)/(pq)$. Since $\Prob(K=k)$ is a sum of path
probabilities, $\partial_q\log\Prob(K=k)=\E[(J-(N-1)q)/(pq)\mid K=k]$. As
$\E J=(N-1)q$, this gives $I_K(q)=\Var(\E[J\mid K])/(pq)^2$, and we divide by
$(N-1)/(pq)=\Var J/(pq)^2$. By the run count of~\cite{paperA}, for $0<k<N$
the probability $\Prob(K=k,J=j)$ is $\frac12p^{N-1}u^j$ times the number of
words with $k$ right steps, $N-k$ left steps and $j$ switches. Summing over
$j$ gives $\Prob(K=k)=\frac12p^{N-1}S_{k,N-k}(u)$, so
$\E[J\mid K=k]=uS'_{k,N-k}(u)/S_{k,N-k}(u)$. Only the two constant paths have
$K\in\{0,N\}$.
\end{proof}

This identity, checked in Lean, says that $e_N$ is the correlation ratio of
$J$ on $K$. Since $\E[J\mid K]$ is the orthogonal projection of $J$ onto
the functions of $K$ in $L^2$, the Cauchy--Schwarz inequality gives
\begin{equation}\label{eq:fisher-proj}
 e_N(q)\ge\frac{\Cov(J,g(K))^2}{\Var J\,\Var g(K)}
 \qquad\text{whenever }\Var g(K)>0.
\end{equation}

\begin{proposition}[few switches]\label{rem:fisher-small}
\begin{enumerate}[(a)]
\item For fixed $N\ge3$, $e_N(q)=1-\frac{N-2}2q+O_N(q^2)$ as $q\to0$.
\item Let $q\to0$ with $T=Nq$ bounded. Then $e_N(q)\to1$ if and only if
$T\to0$.
\end{enumerate}
\end{proposition}

In words, the endpoint loses information as soon as paths with two switches
appear. The reason is that one switch and two switches both leave the endpoint
uniform on the interior bins, so an interior bin cannot tell them apart.

\begin{proof}
By Proposition~\ref{prop:eta}, $1-e_N=\E\Var(J\mid K)/\Var J$, and
$\Var J=(N-1)pq$. Fix an interior bin $0<k<N$. Exactly two words with $k$
right steps have one switch, and exactly $N-2$ have two switches ($k-1$ that
start to the right and $N-k-1$ that start to the left). So
\[
 \Prob(J=1,K=k)=p^{N-2}q,\qquad\Prob(J=2,K=k)=\frac{N-2}2\,p^{N-3}q^2 .
\]

(a) Here $N$ is fixed. Given $K=k$ the number of switches is $1$, or $2$ with
probability $\frac{N-2}2q+O_N(q^2)$, or larger with probability $O_N(q^2)$. So
$\Var(J\mid K=k)=\frac{N-2}2q+O_N(q^2)$, and
$\Prob(0<K<N)=(N-1)q+O_N(q^2)$. For $K\in\{0,N\}$ the conditional variance is
$0$. Hence $\E\Var(J\mid K)=\frac{(N-1)(N-2)}2q^2+O_N(q^3)$, and dividing by
$(N-1)pq$ gives the claim.

(b) For the upper bound we replace $\E[J\mid K]$ by the indicator of an
interior bin, which can only raise the mean squared error. This gives
\[
 \E\Var(J\mid K)\le\E\bigl(J-\mathbf 1\{0<K<N\}\bigr)^2
 =\E\bigl[(J-1)^2;J\ge2\bigr]\le\E[J(J-1)]=(N-1)(N-2)q^2 ,
\]
because $J\ge1$ exactly on the interior bins. So
$1-e_N\le(N-2)q/p\le T/p$, which tends to $0$ if $T\to0$.

For the lower bound, a variable that takes two values $1$ apart with
probabilities $a$ and $b$ has variance at least $ab$. So
$\Var(J\mid K=k)\ge\Prob(J=1\mid K=k)\Prob(J=2\mid K=k)$, and by the
Cauchy--Schwarz inequality
\[
 \E\Var(J\mid K)\ge\sum_{0<k<N}\frac{\Prob(J=1,K=k)\Prob(J=2,K=k)}{\Prob(K=k)}
 \ge\Bigl(\sum_{0<k<N}\sqrt{\Prob(J=1,K=k)\Prob(J=2,K=k)}\Bigr)^2
 =\frac{(N-1)^2(N-2)}2\,p^{2N-5}q^3 .
\]
Dividing by $(N-1)pq$ gives $1-e_N\ge\frac{(N-1)(N-2)}2p^{2N-6}q^2$. If $T$
stays in an interval $[t_0,t_1]$ with $t_0>0$, the right side tends to a limit
between $\frac12t_0^2e^{-2t_1}$ and $\frac12t_1^2$, and it is positive. So
$e_N$ does not tend to $1$ along any sequence on which $T$ stays away from
$0$.
\end{proof}

\begin{theorem}[many switches]\label{thm:fisher}
Let $R^2_N(q)=\Cov(J,X^2)^2/(\Var J\,\Var X^2)$ be the squared correlation
of $J$ and $X^2$. For $N\ge2$ and $0<q<1$,
\[
 R^2_N(q)=\frac1{2(N-1)pq}\cdot
 \frac{\bigl[1+\rho^N-p(1-\rho^N)/T\bigr]^2}{V},
\]
where
\[
 V=1-\frac{1+8\rho+\rho^2+8\rho^{N+1}}{4pT}
 +\frac{\rho(1-\rho^N)\bigl[2(1+2\rho)(2+\rho)-\rho(1-\rho^N)\bigr]}{8p^2T^2}.
\]
There are absolute constants $C$ and $N_0$ such that for $N\ge N_0$,
$0<q\le\frac12$ and $1\le T\le N^{1/4}$,
\[
 0\le e_N(q)-R^2_N(q)\le C\Bigl(\frac1{T^3}+\frac{T^3}N\Bigr).
\]
Consequently, in the same range and with an absolute implied constant,
\[
 e_N(q)=\frac1{2T}+\frac1{4T^2}+O\Bigl(\frac1{T^3}+\frac{T^3}N\Bigr).
\]
\end{theorem}

In words, when $T$ is large but at most $N^{1/4}$, the endpoint keeps about the
fraction $1/(2T)$ of the information, and almost all of what it keeps is carried
by the single statistic $X^2$. The lower bound $e_N\ge R^2_N$ is
\eqref{eq:fisher-proj} with $g(K)=X^2$,
for all $N\ge2$ and $0<q<1$. For the upper bound, Proposition~\ref{prop:eta}
and the uniform Bessel estimate for every bin in~\cite{paperA}
(recalled in~\eqref{eq:fisher-bessel-bound} below) make
$\E[J\mid K]$ close to $y\sqrt{1-W^2}+\frac12$. As $W$ is of order
$T^{-1/2}$, this is close to $y(1-W^2/2)+\frac12$, which lies in the span of
$1$ and $X^2$. Appendix~\ref{app:fisher} proves the closed form and the
bound.

\begin{proof}[Proof of the expansion in Theorem~\ref{thm:fisher}]
If $T\le18$, then $\lvert e_N-\frac1{2T}-\frac1{4T^2}\rvert\le1\le18^3T^{-3}$,
since $e_N$ and $\frac1{2T}+\frac1{4T^2}$ lie in $[0,1]$. Let $T\ge18$. As $q\le\frac12$, we have $p\ge\frac12$
and $0\le\rho<1$, and $\rho^N\le e^{-2Nq}=e^{-2T}\le T^{-2}$. The middle term of $V$ lies
in $[0,9/T]$, and since $1+8\rho+\rho^2=10-20q+4q^2$ it equals
$\frac5{2T}+O(q/T+T^{-2})$. The last term lies in $[0,9/T^2]$. So
$V\ge\frac12$, $1/V=1+\frac5{2T}+O(q/T+T^{-2})$, the squared bracket is
$1-\frac2T+O(q/T+T^{-2})$, and $1/(2(N-1)pq)=(1+O(q+1/N))/(2T)$. Multiplying,
$R^2_N=\frac1{2T}+\frac1{4T^2}+O(T^{-3}+q/T+1/(NT))$ with absolute
constants. Since $q/T=1/N\le T^3/N$, the bound on $e_N-R^2_N$ gives the
expansion.
\end{proof}

\begin{corollary}[at Kagey's crossing]\label{cor:fisher-crossing}
Let $p_N$ be the value of $p$ at which the middle bin and the endpoints are
equally likely, and let $q_N=1-p_N$ and $z_N=Nq_N$
(Section~\ref{sec:setting}). Then
\[
 e_N(q_N)=\frac1{2L}+\frac{\log L+1+\log(8/\pi)}{4L^2}
 +O\Bigl(\frac{(\log L)^2}{L^3}\Bigr).
\]
\end{corollary}

\begin{proof}
Since $S_N(u)\ge2u$, we have $u_N\le\frac12$, so $p_N\ge\frac23$ and
$q_N\le\frac13$. By the second-order expansion of $z_N$ in~\cite{paperA},
$z_N=L(1-\varepsilon)$ with
$\varepsilon=(\frac12\log L-c_0)/L+O(\log L/L^2)$. So Theorem~\ref{thm:fisher}
applies with $T=z_N$ for large $N$ and gives
$e_N(q_N)=\frac1{2z_N}+\frac1{4z_N^2}+O(L^{-3})$. Here
$\frac1{2z_N}=\frac1{2L}+\frac{\log L-2c_0}{4L^2}+O((\log L)^2/L^3)$,
$\frac1{4z_N^2}=\frac1{4L^2}+O(\log L/L^3)$ and $-2c_0=\log(8/\pi)$.
\end{proof}

So $e_N\sim1/(2\log N)$ there, but the relative correction is of order
$(\log\log N)/\log N$ and is large at every $N$ we computed. The product
$2Le_N$ is $1.74$ at $N=100$ and $1.35$ at $N=3200$ (Table~\ref{tab:fisher}).

\begin{table}[ht]
\centering\small
\begin{tabular}{@{}rcccccc@{}}
\toprule
$N$ & $z_N$ & $e_N(q_N)$ & $2z_Ne_N$ & $2Le_N$ & $R^2_N(q_N)$
 & $\frac1{2z_N}+\frac1{4z_N^2}$\\
\midrule
50 & 2.8998 & 0.24117 & 1.3987 & 1.8869 & 0.20873 & 0.20216\\
100 & 3.5034 & 0.18869 & 1.3221 & 1.7379 & 0.16642 & 0.16309\\
400 & 4.7501 & 0.12681 & 1.2047 & 1.5196 & 0.11715 & 0.11634\\
1600 & 6.0246 & 0.09433 & 1.1366 & 1.3919 & 0.09007 & 0.08988\\
3200 & 6.6684 & 0.08361 & 1.1151 & 1.3496 & 0.08070 & 0.08060\\
\bottomrule
\end{tabular}
\caption{The efficiency of the endpoint at Kagey's crossing $q_N=1-p_N$,
computed from the exact law, with $z_N=Nq_N$ and $L=\log N$. The last two
columns are the squared correlation of Theorem~\ref{thm:fisher} and its
expansion to two terms.}
\label{tab:fisher}
\end{table}\subsection*{Many balls}

Now drop $n$ independent balls and record only the histogram of their bins.
How well can $q$ be estimated? By the Cram\'er--Rao bound no unbiased
estimator has variance below $1/(nI_K(q))$, and the maximum likelihood
estimator attains this as $n\to\infty$. The next corollary says that a much
simpler estimator does almost as well. Put $m(q)=\E_qX^2$, a polynomial in
$q$.

\begin{corollary}[the moment estimator]\label{cor:fisher-moment}
Let $X_1,\dots,X_n$ be the endpoints of $n$ independent balls with switching
probability $q$, and suppose that $R^2_N(q)>0$. Then $m'(q)\ne0$, and with
probability tending to $1$ the equation
$m(\hat q)=\frac1n\sum_{i=1}^nX_i^2$ has a solution $\hat q_n$ that converges
to $q$ in probability. For this solution
\[
 \sqrt n\,(\hat q_n-q)\Rightarrow N\Bigl(0,\frac{pq}{(N-1)R^2_N(q)}\Bigr)
 \qquad(n\to\infty).
\]
So its asymptotic efficiency, relative to the bound $1/(nI_K(q))$, is
$R^2_N(q)/e_N(q)$. There is an absolute constant $C'$ such that, with $N_0$
as in Theorem~\ref{thm:fisher}, for $N\ge N_0$, $0<q\le\frac12$ and
$1\le T\le N^{1/4}$,
\[
 1-C'\Bigl(\frac1{T^2}+\frac{T^4}N\Bigr)\le\frac{R^2_N(q)}{e_N(q)}\le1 .
\]
\end{corollary}

In words, among all estimators that see only the bins, matching the mean
square of the endpoint loses a fraction of order $1/T^2+T^4/N$ of the
information. At Kagey's crossing the efficiency $R^2_N/e_N$ is $0.87$ at
$N=50$ and $0.97$ at $N=3200$ (Table~\ref{tab:fisher}).

\begin{proof}
The law of a path is a polynomial in $q$ with score $(J-(N-1)q)/(pq)$, so
$m'(q)=\E[X^2(J-(N-1)q)]/(pq)=\Cov(J,X^2)/(pq)$. This is not zero because
$R^2_N(q)>0$. So $m$ is strictly monotone near $q$. By the law of large
numbers the mean of the $X_i^2$ tends to $m(q)$, so with probability tending
to $1$ it lies in the image of a small interval around $q$, and the solution
in that interval converges to $q$. The central limit theorem and the delta
method give the limit law with variance
\[
 \frac{\Var X^2}{m'(q)^2}=\frac{(pq)^2\Var X^2}{\Cov(J,X^2)^2}
 =\frac{(pq)^2}{R^2_N(q)\Var J}=\frac{pq}{(N-1)R^2_N(q)} .
\]
Since $I_K(q)=e_N(q)(N-1)/(pq)$, the ratio of $1/I_K(q)$ to this variance is
$R^2_N/e_N$, and it is at most $1$ by~\eqref{eq:fisher-proj}. For the lower
bound let $C_1$ be the constant in the expansion of
Theorem~\ref{thm:fisher}, and put $C'=4\max(C,C_1)$ and
$\epsilon=T^{-2}+T^4/N$. If $C'\epsilon\ge1$ the claim is trivial, since
$R^2_N/e_N\ge0$. Otherwise $C_1\epsilon<\frac14$, so the expansion gives
$e_N\ge\frac1{2T}-\frac{C_1\epsilon}T\ge\frac1{4T}$. Then
$1-R^2_N/e_N=(e_N-R^2_N)/e_N\le4T\cdot C\epsilon/T=4C\epsilon\le C'\epsilon$.
\end{proof}
